%%%%%%%%%%%%%%%%%%%%%%%%%%%%%%%%%%%%%%%%%%%%%%%%%%%%%%%%%%%%%%%%%%%
%                                                                 %
%  GEANT manual in LaTeX form                                     %
%                                                                 %
%  Michel Goossens (for translation into LaTeX)                   %
%  Version 1.00                                                   %
%  Last Mod. Jan 24 1991  1300   MG + IB                          %
%                                                                 %
%%%%%%%%%%%%%%%%%%%%%%%%%%%%%%%%%%%%%%%%%%%%%%%%%%%%%%%%%%%%%%%%%%%
\Origin{P.Zanarini}   
\Revision{S.Giani}
\Documentation{P.Zanarini, S.Giani}
\Submitted{15.01.84}            \Revised{10.12.93}
\Version{Geant 3.16}            \Routid{DRAW210}

\Makehead{Drawing the geometrical tree}
The geometrical tree is a hierarchical representation of the 
structure of the detector, based on the mother-daughter
relationship among the various volumes composing it. The \Rind{GDTREE} 
routine (corresponding to the {\tt DTREE} interactive command) allows 
the drawing of the logical tree, displaying name, multiplicity 
and other information about the volumes.

\Shubr{GDTREE}{(CHNAME,LEVMAX,ISEL)}
\begin{DLtt}{MMMMMMMM}
\item[CHNAME]  ({\tt CHARACTER*4}) volume name;
\item[LEVMAX]  ({\tt INTEGER}) depth levels (if 0 all levels are taken) 
default is 3;
\item[ISEL]    ({\tt INTEGER}) used to select options in the picture of the tree:
\begin{DLtt}{MMMMMM}
\item[~~~~0] only the node name is drawn;
\item[xxxx1] multiplicity is added: 
\begin{itemize}
\item each node that has been positioned several 
times (by \Rind{GSPOS} or \Rind{GSPOSP}) has a multiplicity number preceded by 
the symbol '*';
\item each node that has been obtained by division mechanism
(by \Rind{GSDVN}, \Rind{GSDVT}, etc.) has a multiplicity number preceded by the
symbol \#;
\end{itemize}
\item[xxx1x] {\tt ONLY} information is added.
volumes that are {\tt MANY} will be
drawn as three square nodes, one overlapping the others but slightly shifted;

{\bf Note:} if {\tt CHNAME} itself is {\tt MANY} it will not be recognised as 
such, because this information is stored in its mother which is unknown;

\item[xx1xx] {\tt DET} information is added; if a volumes is a detector it 
is drawn as a square node with another one slightly larger around it;
\item[x1xxx] {\tt SEEN} information is added;
if a volume is not seen (either because of its {\tt SEEN} attribute or
because of its ancestors {\tt SEEN} attributes affecting the whole tree path)
its square node is drawn by using dashed lines;
\item[1xxxx] a little picture of the volume is added above each node.
\end{DLtt}
\end{DLtt}

\begin{figure}[hbt]
     \centering
     \epsfig{file=eps/draw210-1.eps,width=14cm}
\begin{verbatim}
      CALL GDTREE('FD  ',5,11111)
\end{verbatim}
     \caption{Example of geometrical tree structure}
     \label{fg:draw210-1}
\end{figure}
\Rind{GDTREE} draws the tree of the volume data structure starting
from the volume {\tt CHNAME}, for {\tt LEVMAX} depth levels;
a continuation link line is drawn for those elements
at the last level on the picture but having descendants.
An example of use of \Rind{GDTREE} can be found in Fig.~\ref{fg:draw210-1}.
For every node, as many red arrows as the number of existing levels below it
are plotted on the link; vice versa, as many green arrows as the number of 
existing levels before it are also plotted.
When running {\tt GEANT} interactively, if the {\tt ISEL}
argument is not specified, control is given to the
mouse after issuing the {\tt DTREE} command. This allows the display of various
volume parameters and the navigation of the geometrical tree by
clicking the mouse buttons.
The size of the picture and the dimensions of each
element and arrows are computed automatically. 

\section*{Interactive use}
In an interactive session, control can be given to the mouse: 
\begin{itemize}
\item clicking on the left button when the cursor is inside a volume
{\it pave} will generate a \Rind{DSPEC} command for that volume; 
\item clicking the same button when the cursor is on a red arrow will 
perform a \Rind{DTREE} for the volume preceding the arrow, displaying a
number of levels which depends on the which arrow has been selected;
\item selecting in the same way the i$^{th}$ green arrow preceding 
a given volume will perform a \Rind{DTREE}
for its mother-volume staying {\it i} levels before;
\item clicking the right button of the mouse will return
the control to the command mode.
\end{itemize}
 
When running with the X-windows software, the drawing of the 
specification (\Rind{DSPEC}) appears in a different window to 
speed up the scanning of the tree. Iterating this procedure it is 
possible to analyse very easily and quickly any kind of geometry. 
This drawing tree capability can be very
useful when tuning the detector geometry.
See the interactive section {\tt [XINT]} for examples.
